\documentclass[11pt]{article}

% ---------- Packages ----------
\usepackage[a4paper, margin=1in]{geometry}
\usepackage{graphicx}
\usepackage{booktabs}
\usepackage{longtable}
\usepackage{array}
\usepackage[table]{xcolor}
\usepackage{titlesec}
\usepackage{enumitem}
\usepackage{caption}
\usepackage{float}
\usepackage{fancyhdr}
\usepackage{hyperref}
\usepackage{amsmath}
\usepackage{amssymb}
\usepackage{parskip}
\usepackage{tabularx}

% ---------- Colors ----------
\definecolor{headingblue}{RGB}{31,56,100}
\definecolor{subblue}{RGB}{46,83,149}
\definecolor{tableheadbg}{RGB}{46,83,149}
\definecolor{lightgray}{RGB}{136,136,136}

% ---------- Heading styles ----------
\titleformat{\section}{\normalfont\Large\bfseries\color{headingblue}}{\thesection}{0.6em}{}
\titleformat{\subsection}{\normalfont\large\bfseries\color{subblue}}{\thesubsection}{0.6em}{}
\titleformat{\subsubsection}{\normalfont\normalsize\bfseries\color{black}}{\thesubsubsection}{0.6em}{}
\titlespacing*{\section}{0pt}{18pt}{10pt}
\titlespacing*{\subsection}{0pt}{14pt}{8pt}
\titlespacing*{\subsubsection}{0pt}{10pt}{6pt}

% ---------- Header/Footer ----------
\pagestyle{fancy}
\fancyhf{}
\renewcommand{\headrulewidth}{0pt}
\fancyhead[R]{\small\color{lightgray}Student Placement Prediction --- Case Study}
\fancyfoot[C]{\thepage}

% ---------- Hyperref ----------
\hypersetup{
    colorlinks=true,
    linkcolor=subblue,
    citecolor=subblue,
    urlcolor=subblue,
    pdftitle={Student Placement Prediction: A Data Analysis Case Study}
}

% ---------- Helper commands ----------
\newcommand{\fig}[3]{%
    \begin{figure}[H]
        \centering
        \includegraphics[width=#2\linewidth]{images/#1}
        \caption{#3}
    \end{figure}
}

\newcommand{\reject}{\textbf{Reject } $H_0$}
\newcommand{\failreject}{\textbf{Fail to reject } $H_0$}

\begin{document}

% ======================================================
% TITLE PAGE
% ======================================================
\begin{titlepage}
\centering
\vspace*{2.5cm}

{\large\textsc{Case Study}}\\[1cm]

{\huge\bfseries\color{headingblue} Student Placement Prediction:}\\[0.3cm]
{\huge\bfseries\color{headingblue} A Data Analysis Case Study}\\[1.2cm]

{\itshape\large\color{gray} An Exploratory Data Analysis and Statistical Study of\\
Student Demographic, Academic, and Placement Data}\\[2.5cm]

\begin{tabularx}{0.75\linewidth}{|>{\bfseries}p{4cm}|X|}
\hline
Team Members & \rule{0pt}{1em}\underline{\hspace{6cm}} \\ \hline
Roll Number(s) & \underline{\hspace{6cm}} \\ \hline
Course / Subject & \underline{\hspace{6cm}} \\ \hline
Institution & \underline{\hspace{6cm}} \\ \hline
Guided By & \underline{\hspace{6cm}} \\ \hline
Date of Submission & \underline{\hspace{6cm}} \\ \hline
\end{tabularx}

\vfill
\end{titlepage}

% ======================================================
% TABLE OF CONTENTS
% ======================================================
\tableofcontents
\newpage

% ======================================================
% 2. PROBLEM STATEMENT
% ======================================================
\section{Problem Statement}

Campus placement outcomes are of central concern to students, academic institutions, and recruiters alike. Institutions want to understand which student attributes --- demographic (age, gender, city), academic (CGPA, degree/profession level), or otherwise --- are meaningfully associated with whether a student is placed in a job after graduation. Understanding these relationships helps institutions design better placement support programs, and helps students understand what factors genuinely affect employability outcomes versus common assumptions and myths.

This case study analyzes a dataset of 1{,}100 student records containing demographic and academic attributes alongside a binary placement outcome. The objective is to clean and prepare the raw data, explore it visually, and apply formal statistical tests to answer the central question:

\begin{center}
\itshape ``Which student attributes, if any, are statistically associated with placement status?''
\end{center}

The analysis follows a complete data-science workflow --- data cleaning, transformation, encoding, outlier handling, visualization, and hypothesis testing --- culminating in evidence-based conclusions about the drivers (or absence thereof) of student placement in this dataset.

% ======================================================
% 3. INFORMATION OF THE DATA
% ======================================================
\section{Information of the Data}

The raw dataset contains 1{,}100 rows and 7 columns, mixing categorical (nominal) and numerical (continuous) features, along with a binary target variable. Table~\ref{tab:columns} summarizes each column.

\begin{table}[H]
\centering
\caption{Dataset column summary}
\label{tab:columns}
\begin{tabularx}{\linewidth}{l l X l}
\toprule
\rowcolor{tableheadbg}
\textcolor{white}{\textbf{Column}} & \textcolor{white}{\textbf{Type}} & \textcolor{white}{\textbf{Description}} & \textcolor{white}{\textbf{Missing (raw)}} \\
\midrule
\texttt{name} & Categorical (text) & Student's first name & 85 (7.7\%) \\
\texttt{city} & Categorical (nominal) & City of residence (asgard, wakanda, gotham, purgatory) & 184 (16.7\%) \\
\texttt{gender} & Categorical (binary) & Gender of the student (male / female) & 99 (9.0\%) \\
\texttt{profession} & Categorical (ordinal) & Highest degree pursued (bachelor / masters / phd) & 169 (15.4\%) \\
\texttt{age} & Numeric (continuous) & Student's age in years & 112 (10.2\%) \\
\texttt{cgpa} & Numeric (continuous) & Cumulative Grade Point Average (0--10 scale) & 132 (12.0\%) \\
\texttt{placed} & Binary (target) & Placement outcome (1 = placed, 0 = not placed) & 0 (0\%) \\
\bottomrule
\end{tabularx}
\end{table}

The dataset shape was verified using \texttt{df.shape} and \texttt{df.info()}, confirming 1{,}100 entries across 7 columns with the null counts and dtypes shown below.

\fig{data_info.png}{0.75}{Dataset shape, sample rows, and \texttt{df.info()} output}

The target variable \texttt{placed} is well balanced, with 521 students not placed (50.3\%) and 514 students placed (49.7\%) after cleaning --- an almost even 50/50 split, which is favorable for unbiased statistical comparison between the two groups.

% ======================================================
% 4. SOURCE OF DATA
% ======================================================
\section{Source of Data}

The dataset (\texttt{my\_data.csv}) was provided as a course/assignment dataset for the purpose of this case study. It is a structured, tabular dataset resembling a typical student-placement record used in academic exercises on data cleaning and statistical analysis.

It is worth noting that the categorical values in the \texttt{name} column (e.g., moriarity, holmes, sherlock, sam, dean, castiel, bobby) and in the \texttt{city} column (asgard, wakanda, gotham, purgatory) are drawn from popular fictional universes rather than real-world identifiers. This strongly suggests the dataset is synthetically generated / anonymized for educational purposes rather than sourced from a real institutional placement cell. This context is relevant when interpreting results --- patterns found here characterize the dataset itself and the methodology applied to it, not any real population of students.

% ======================================================
% 5. DATA PREPROCESSING
% ======================================================
\section{Data Preprocessing}

\subsection{a. Data Cleaning}

\subsubsection*{Duplicate Records}
Duplicate rows were checked with \texttt{df.duplicated().sum()}. The raw dataset (1{,}100 $\times$ 7) contained 65 fully duplicate rows (5.9\% of the data), which were removed using \texttt{df.drop\_duplicates()}. A later verification pass, run after feature engineering (standardization, binning, and encoding had expanded the frame to 16 columns), confirmed 1{,}032 unique rows with zero remaining duplicates.

\fig{duplicate_data.jpg}{0.75}{Detecting the 65 duplicate rows before removal}
\fig{drop_duplicate.jpg}{0.75}{Shape confirmation after dropping duplicates}

\subsubsection*{Missing Value Detection and Treatment}
All seven columns except the target were found to contain missing values, ranging from 8.2\% (\texttt{name}) to 17.8\% (\texttt{city}). Given the scale of missingness, rows were not dropped outright (this would have discarded a large fraction of the dataset); instead, column-appropriate imputation was applied:

\begin{itemize}[leftmargin=1.2cm]
    \item \textbf{Numerical columns} (\texttt{age}, \texttt{cgpa}): imputed with the column \textbf{median}, which is robust to skew and outliers.
    \item \textbf{Categorical columns} (\texttt{city}, \texttt{gender}, \texttt{profession}): imputed with the column \textbf{mode} (most frequent category).
    \item \texttt{name}: imputed with the placeholder label ``Unknown'', since it carries no analytical/statistical value beyond identification.
\end{itemize}

\fig{handling_missing_data.jpg}{0.8}{Missing value counts and percentage-missing per column}
\fig{handling_age_cgpa.jpg}{0.7}{Median imputation for \texttt{age} and \texttt{cgpa}}
\fig{handling_name_.jpg}{0.7}{Mode imputation for \texttt{city}/\texttt{gender}/\texttt{profession} and placeholder fill for \texttt{name}}

After imputation, \texttt{df.isnull().sum()} confirmed zero missing values across all columns.

\subsubsection*{Data Type Correction}
\texttt{df.dtypes} revealed that \texttt{age} was stored as \texttt{float64} despite representing whole years. It was cast to an integer type using \texttt{df["age"] = df["age"].astype(int)} for consistency and readability.

\fig{datatype_in_column.jpg}{0.7}{Correcting the dtype of the \texttt{age} column}

\subsection{b. Data Standardization}

The numerical features \texttt{age} and \texttt{cgpa} were standardized using scikit-learn's \texttt{StandardScaler}, which rescales values to zero mean and unit variance ($z$-scores). This step is essential before running distance- or scale-sensitive analyses, and was verified by re-computing the mean ($\approx 0$) and standard deviation ($=1.0$) of the resulting \texttt{age\_standardized} and \texttt{cgpa\_standardized} columns.

\fig{standardization.jpg}{0.65}{\texttt{StandardScaler} applied to \texttt{age} and \texttt{cgpa}, with verification statistics}

For reference, the descriptive statistics of the cleaned (pre-standardization) numeric columns were:

\fig{stats_before_standarization.jpg}{0.55}{\texttt{df.describe()} after cleaning, before standardization}

\subsection{c. Data Transformation: Binning}

To support categorical-style analysis (e.g., count plots, chi-square tests) alongside the continuous variables, two new binned features were engineered using \texttt{pd.cut()}:

\begin{itemize}[leftmargin=1.2cm]
    \item \textbf{Age\_Group}: Young (18--25), Adult (25--35), Senior (35--60).
    \item \textbf{CGPA\_Category}: Low (4--6), Average (6--8), High (8--10).
\end{itemize}

\fig{tansformation_01.jpg}{0.8}{Binning code for \texttt{Age\_Group} and \texttt{CGPA\_Category}}
\fig{tansformation_02.jpg}{0.55}{Resulting distribution of \texttt{Age\_Group} and \texttt{CGPA\_Category}}

The resulting distribution is fairly balanced: 437 Young, 398 Adult, and 200 Senior students; and 423 High, 412 Average, and 200 Low CGPA students.

\subsection{d. Handling Categorical Data: Encoding Techniques}

Before encoding, a readable copy of the dataframe (\texttt{df\_viz}) was preserved with the original human-readable category labels, specifically for use in visualizations and hypothesis testing where readability matters more than numeric encoding.

\fig{encoding.jpg}{0.85}{Preserving a readable copy (\texttt{df\_viz}) before encoding}

\texttt{gender}, a binary categorical variable, was transformed using \textbf{Label Encoding} (scikit-learn's \texttt{LabelEncoder}), mapping female $\rightarrow$ 0 and male $\rightarrow$ 1.

\fig{labelencoding.jpg}{0.85}{Label encoding applied to \texttt{gender}}

\texttt{city} and \texttt{profession}, both nominal categorical variables with more than two categories, were transformed using \textbf{One-Hot Encoding} (\texttt{pd.get\_dummies}), avoiding the false ordinal relationship that label encoding would have implied for these columns.

\fig{one_hot_encoding.jpg}{0.85}{One-hot encoding of \texttt{city} and \texttt{profession}}

\subsection{e. Outlier Detection}

Outliers in the two continuous variables (\texttt{age}, \texttt{cgpa}) were detected using two complementary methods: visual inspection via box plots, and the mathematical Interquartile Range (IQR) rule (values beyond $Q1 - 1.5\times IQR$ or $Q3 + 1.5\times IQR$).

\fig{outliers_age_b.jpg}{0.55}{Box plot of \texttt{age}, showing high-end outliers above $\sim$42 years}
\fig{outliers_age_i.jpg}{0.5}{IQR bounds for \texttt{age} (10.0, 42.0) and the 154 rows identified as outliers}
\fig{outliers_cgpa_b.jpg}{0.5}{Box plot of CGPA, showing a cluster of low-end outliers below $\sim$5.2}

Rather than deleting outlier rows outright (which would have discarded a further $\sim$15\% of the data), extreme values were clipped to the IQR boundary limits using \texttt{df["age"] = df["age"].clip(lower\_limit, upper\_limit)}. A follow-up check confirmed zero rows remained outside the bounds after clipping.

\fig{handling_outliers.jpg}{0.7}{Clipping outliers to IQR boundaries and confirming zero out-of-bound rows}

After all cleaning and outlier handling, the final numeric summary (age: mean 27.74, range 18--42; cgpa: mean 7.61, range 4.01--9.49; placed: mean 0.497) reflects a clean, well-behaved dataset of 1{,}035 rows ready for analysis.

% ======================================================
% 6. DATA VISUALIZATION
% ======================================================
\section{Data Visualization}

A range of univariate and bivariate plots were generated using \texttt{matplotlib} and \texttt{seaborn} to explore the distribution of each feature and its relationship with the placement outcome.

\subsection{Univariate Distributions}

\fig{age_cgpa_distribution.jpg}{0.5}{Distribution of Age (mildly bimodal, with a secondary cluster around 40--42) and CGPA (right-skewed, with a secondary low-scoring cluster below 6)}
\fig{gender_CP.jpg}{0.55}{Gender distribution: males ($\sim$813) considerably outnumber females ($\sim$220) in the dataset}
\fig{city_CP.jpg}{0.6}{City distribution: Wakanda is the dominant city ($\sim$547 students), followed by Gotham, Asgard, and Purgatory}
\fig{profession_CP.jpg}{0.55}{Profession (degree level) distribution: Bachelor's students dominate ($\sim$657), followed by Master's ($\sim$254) and PhD ($\sim$119)}

The target variable itself is balanced: 521 students not placed vs. 514 placed (50.3\% vs. 49.7\%), which avoids the class-imbalance issues that can complicate group comparisons.

\subsection{Placement vs. Categorical Features}

\fig{genderVSplacement.png}{0.55}{Placement counts by Gender}
\fig{cityVSplacement.png}{0.6}{Placement counts by City}
\fig{placementVSprofession.png}{0.55}{Placement counts by Profession}
\fig{age_groupVSplacement.png}{0.55}{Placement counts by Age Group}
\fig{cgpa_catogaryVSplacement.png}{0.55}{Placement counts by CGPA Category}

The placement rate (\% placed) within each category, computed with \texttt{groupby().mean()}, was:

\begin{table}[H]
\centering
\caption{Placement rate by Gender}
\begin{tabular}{ll}
\toprule
\rowcolor{tableheadbg}
\textcolor{white}{\textbf{Gender}} & \textcolor{white}{\textbf{Placement Rate}} \\
\midrule
Female & 42.99\% \\
Male & 51.47\% \\
\bottomrule
\end{tabular}
\end{table}

\begin{table}[H]
\centering
\caption{Placement rate by Profession}
\begin{tabular}{ll}
\toprule
\rowcolor{tableheadbg}
\textcolor{white}{\textbf{Profession}} & \textcolor{white}{\textbf{Placement Rate}} \\
\midrule
Bachelor & 50.08\% \\
Masters & 48.44\% \\
PhD & 50.00\% \\
\bottomrule
\end{tabular}
\end{table}

\begin{table}[H]
\centering
\caption{Placement rate by City}
\begin{tabular}{ll}
\toprule
\rowcolor{tableheadbg}
\textcolor{white}{\textbf{City}} & \textcolor{white}{\textbf{Placement Rate}} \\
\midrule
Asgard & 46.15\% \\
Gotham & 52.36\% \\
Purgatory & 45.95\% \\
Wakanda & 50.18\% \\
\bottomrule
\end{tabular}
\end{table}

Gender shows the widest gap in raw placement rate (male 51.5\% vs. female 43.0\%) among the categorical breakdowns, foreshadowing the statistically significant chi-square result found in Section~7.

\subsection{Placement vs. Numeric Features}

\fig{ageVSgpa_scatterplot.jpg}{0.75}{Age vs. CGPA scatter plot, colored by placement status --- no visible separation between the two placement groups}
\fig{ageVSplaced_BP.png}{0.5}{Age distribution by placement status (box plot) --- near-identical medians and spread}
\fig{cgpaVSplaced_BP.png}{0.5}{CGPA distribution by placement status (box plot) --- near-identical medians and spread}

The descriptive statistics confirm the near-total overlap visible in the box plots:

\fig{stats_for_BP.jpg}{0.75}{Age and CGPA summary statistics grouped by placement status}

\subsection{Correlation Heatmap}

\fig{heatmap.png}{0.6}{Correlation heatmap of \texttt{age}, \texttt{cgpa}, \texttt{gender}, and \texttt{placed}}

All pairwise correlations among \texttt{age}, \texttt{cgpa}, \texttt{gender}, and \texttt{placed} are very close to zero (largest magnitude: 0.07 between \texttt{gender} and \texttt{placed}), visually confirming the absence of any strong linear relationship among these variables.

% ======================================================
% 7. STATISTICAL ANALYSIS
% ======================================================
\section{Statistical Analysis}

Beyond visual inspection, formal statistical hypothesis tests were used to determine whether observed differences and associations in the data are statistically significant ($\alpha = 0.05$) or could plausibly be due to random chance.

\subsection{a. Hypothesis Testing}

\subsubsection*{Test 1: Independent $t$-test --- CGPA vs. Placed}
\textit{$H_0$: There is no significant difference in mean CGPA between placed and not-placed students.}\\
\textit{$H_1$: There is a significant difference in mean CGPA between placed and not-placed students.}

\fig{independent_t_test_cgpaVSplaced.jpg}{0.85}{Independent $t$-test, CGPA vs. Placed: $t = 0.1197$, $p = 0.9048$}

\textbf{Result:} $p = 0.9048 > 0.05 \Rightarrow$ \failreject. No statistically significant difference in CGPA between placed and not-placed students.

\subsubsection*{Test 2: Independent $t$-test --- Age vs. Placed}
\textit{$H_0$: There is no significant difference in mean Age between placed and not-placed students.}\\
\textit{$H_1$: There is a significant difference in mean Age between placed and not-placed students.}

\fig{independent_t_test_ageVSplaced.jpg}{0.85}{Independent $t$-test, Age vs. Placed: $t = 0.2921$, $p = 0.7703$}

\textbf{Result:} $p = 0.7703 > 0.05 \Rightarrow$ \failreject. No statistically significant difference in Age between placed and not-placed students.

\subsubsection*{Test 3: Chi-square Test --- Gender vs. Placed}
\textit{$H_0$: Gender and Placement are independent (no association).}\\
\textit{$H_1$: Gender and Placement are dependent (there is an association).}

\fig{chisquare_test__genderVSplaced.jpg}{0.85}{Chi-square test, Gender vs. Placed: $\chi^2 = 4.6751$, $df = 1$, $p = 0.0306$}

\textbf{Result:} $p = 0.0306 < 0.05 \Rightarrow$ \reject. There is a statistically significant association between Gender and Placement.

\subsubsection*{Test 4: Chi-square Test --- Profession vs. Placed}
\textit{$H_0$: Profession and Placement are independent (no association).}\\
\textit{$H_1$: Profession and Placement are dependent (there is an association).}

\fig{chisquare_test_professionVSplaced.jpg}{0.85}{Chi-square test, Profession vs. Placed: $\chi^2 = 0.2042$, $df = 2$, $p = 0.9029$}

\textbf{Result:} $p = 0.9029 > 0.05 \Rightarrow$ \failreject. No statistically significant association between Profession and Placement.

\subsubsection*{Summary of Hypothesis Tests}

\begin{table}[H]
\centering
\caption{Summary of hypothesis test results}
\begin{tabularx}{\linewidth}{X l l l}
\toprule
\rowcolor{tableheadbg}
\textcolor{white}{\textbf{Test}} & \textcolor{white}{\textbf{Statistic}} & \textcolor{white}{\textbf{$p$-value}} & \textcolor{white}{\textbf{Conclusion ($\alpha=0.05$)}} \\
\midrule
CGPA vs Placed ($t$-test) & $t=0.1197$ & 0.9048 & Not significant \\
Age vs Placed ($t$-test) & $t=0.2921$ & 0.7703 & Not significant \\
Gender vs Placed ($\chi^2$) & $\chi^2=4.6751$ & 0.0306 & Significant \\
Profession vs Placed ($\chi^2$) & $\chi^2=0.2042$ & 0.9029 & Not significant \\
\bottomrule
\end{tabularx}
\end{table}

\subsection{b. Correlation Analysis}

Pearson correlation coefficients were computed for each pair of numeric variables, along with their statistical significance:

\fig{corelation_Analysis.jpg}{0.85}{Pearson correlation coefficients and full correlation matrix}

\begin{table}[H]
\centering
\caption{Pearson correlation summary}
\begin{tabularx}{\linewidth}{X l l l}
\toprule
\rowcolor{tableheadbg}
\textcolor{white}{\textbf{Pair}} & \textcolor{white}{\textbf{Pearson $r$}} & \textcolor{white}{\textbf{$p$-value}} & \textcolor{white}{\textbf{Significant?}} \\
\midrule
Age vs CGPA & $-0.0404$ & 0.1945 & No \\
Age vs Placed & 0.0091 & 0.7703 & No \\
CGPA vs Placed & 0.0037 & 0.9048 & No \\
Gender vs Placed & 0.070 & --- & Consistent w/ $\chi^2$ result \\
\bottomrule
\end{tabularx}
\end{table}

Every correlation coefficient is close to zero ($|r| < 0.08$), confirming that none of the numeric variables share a meaningful linear relationship with each other or with the placement outcome, consistent with the heatmap in Figure~11 and the hypothesis test results above.

% ======================================================
% 8. RESULTS AND DISCUSSION
% ======================================================
\section{Results and Discussion}

Combining the visual exploration and formal statistical testing, the following conclusions emerge from this dataset:

\begin{itemize}[leftmargin=1.2cm]
    \item \textbf{Gender} is the only feature with a statistically significant association with placement ($\chi^2 = 4.6751$, $p = 0.0306$). Male students show a placement rate of 51.5\% versus 43.0\% for female students --- a gap of roughly 8.5 percentage points.
    \item \textbf{CGPA} shows no significant relationship with placement, either via $t$-test ($p = 0.9048$) or correlation ($r = 0.0037$, $p = 0.9048$). Students who were placed and not placed have virtually identical CGPA distributions (means of 7.61 vs. 7.60).
    \item \textbf{Age} shows no significant relationship with placement, either via $t$-test ($p = 0.7703$) or correlation ($r = 0.0091$, $p = 0.7703$).
    \item \textbf{Profession} (degree level) shows no significant association with placement ($\chi^2 = 0.2042$, $p = 0.9029$); bachelor's, master's, and PhD students are placed at almost identical rates ($\sim$48--50\%).
    \item \textbf{City} shows a modest spread in raw placement rates (45.9\% to 52.4\%) but was not formally tested; given the pattern seen with profession, this spread likely reflects sampling noise rather than a true effect.
    \item \textbf{Age and CGPA} are themselves uncorrelated ($r = -0.0404$, not significant), indicating these two numeric features vary independently of one another in this dataset.
\end{itemize}

Overall, the data suggests that in this dataset, placement outcome behaves close to a random 50/50 split with respect to academic performance (CGPA), age, and degree level --- the only detectable (and still modest) signal is a gender-based difference. This is a somewhat counter-intuitive finding relative to common assumptions (e.g., that CGPA strongly predicts placement), and is worth flagging explicitly rather than assuming CGPA ``must'' matter.

Given the earlier observation that the dataset's names and cities are drawn from fictional universes, it is plausible that placement outcomes here were generated with a random or near-random process (possibly with a small deliberate gender effect), rather than reflecting a genuine underlying real-world mechanism. This context should be kept in mind: the statistical conclusions are valid descriptions of this dataset, but should not be over-generalized to real placement processes without validation on real institutional data.

% ======================================================
% 9. CHALLENGES AND FUTURE ENHANCEMENT
% ======================================================
\section{Challenges and Future Enhancement}

\subsection*{Challenges Encountered}
\begin{itemize}[leftmargin=1.2cm]
    \item High missing-value rates (8\%--18\%) across nearly every column required careful, column-specific imputation strategy rather than simple row deletion, to avoid losing a large share of the dataset.
    \item A meaningful number of exact duplicate rows (5.9\%) had to be identified and removed to avoid double-counting records in downstream statistics.
    \item Right- and left-skewed distributions (age's bimodal shape; CGPA's low-end cluster) required outlier handling (IQR clipping) rather than naive $z$-score filtering, to avoid discarding legitimate data.
    \item The weak-to-negligible correlations and mostly non-significant test results meant the dataset offered limited explanatory power for placement --- a reminder that clean methodology can still yield a ``no strong effect'' conclusion, which is itself a valid and useful finding.
    \item Categorical encoding choices (label vs. one-hot) had to be made carefully to avoid introducing false ordinal relationships into nominal variables like \texttt{city}.
\end{itemize}

\subsection*{Future Enhancements}
\begin{itemize}[leftmargin=1.2cm]
    \item Incorporate richer features that plausibly drive placement in the real world --- interview performance, technical/soft-skill assessments, internship experience, extracurricular involvement, and communication skills.
    \item Apply predictive machine learning models (e.g., Logistic Regression, Random Forest, Gradient Boosting) to formally test predictive power and quantify feature importance, going beyond pairwise hypothesis tests to a multivariate model.
    \item Use multivariate/interaction analysis (e.g., two-way ANOVA, interaction terms) to check whether combinations of features (e.g., gender $\times$ profession) reveal effects not visible in single-variable tests.
    \item Validate findings on a real, verified institutional placement dataset to confirm whether the gender effect observed here generalizes, and to check it does not stem from confounding or sampling artifacts.
    \item Apply more advanced missing-data techniques (e.g., KNN or iterative/MICE imputation) instead of simple median/mode imputation, to better preserve relationships between correlated features when imputing.
    \item Build an interactive dashboard (e.g., in Power BI, Tableau, or a Streamlit app) to allow stakeholders to explore placement trends dynamically across cities, professions, and demographic segments.
\end{itemize}

\end{document}
